\documentclass[11pt,a4paper]{article}

\usepackage[utf8]{inputenc}
\usepackage[T1]{fontenc}
\usepackage{lmodern}
\usepackage[margin=2.5cm]{geometry}
\usepackage{booktabs}
\usepackage{tabularx}
\usepackage{array}
\usepackage{enumitem}
\usepackage{xcolor}
\usepackage{titlesec}
\usepackage{fancyhdr}
\usepackage[hidelinks]{hyperref}

\definecolor{accent}{HTML}{1D4ED8}
\titleformat{\section}{\Large\bfseries\color{accent}}{\thesection}{0.8em}{}
\titleformat{\subsection}{\large\bfseries}{\thesubsection}{0.8em}{}
\setlist{itemsep=2pt, topsep=4pt}
\setlength{\parskip}{6pt}
\setlength{\parindent}{0pt}
\renewcommand{\arraystretch}{1.25}
\newcolumntype{L}{>{\raggedright\arraybackslash}X}

\pagestyle{fancy}
\fancyhf{}
\lhead{ditto-chaos: Project Report}
\rhead{\thepage}
\renewcommand{\headrulewidth}{0.4pt}

% ---- Fill these in before submitting -------------------------------------
\newcommand{\teamname}{[Team name]}
\newcommand{\teammembers}{[Member 1], [Member 2], [Member 3], [Member 4]}
\newcommand{\institution}{[Institution / event name]}
\newcommand{\reportdate}{October 2026}
% --------------------------------------------------------------------------

\begin{document}

\begin{titlepage}
  \centering
  \vspace*{3cm}
  {\Huge\bfseries\color{accent} ditto-chaos\par}
  \vspace{0.6cm}
  {\Large Safety-First Chaos Engineering for Digital Twins\par}
  \vspace{2cm}
  {\large Project Report\par}
  \vspace{2.5cm}
  {\large \teamname\par}
  \vspace{0.3cm}
  {\teammembers\par}
  \vspace{1.5cm}
  {\institution\par}
  \vspace{0.3cm}
  {\reportdate\par}
  \vfill
  {\small Source code: \url{https://github.com/NJ800/ditto-chaos}\par}
\end{titlepage}

\begin{abstract}
\noindent
ditto-chaos is a chaos-engineering tool for Eclipse Ditto, an open-source
digital-twin platform. It injects one controlled fault at a time into a local
Ditto sandbox, measures the effect with the same signals an on-call engineer
would use, and stops and rolls back automatically when a safety limit is
crossed. Every rollback is verified by inspecting the system, not assumed from
a command's exit code. After each experiment the tool produces findings and
recommended fixes. Artificial intelligence is optional and never makes a safety
decision: deterministic code authorizes, stops and reverts every experiment.
In a live run, six fault types were injected in sequence; all six were rolled
back and verified, and Ditto recovered in 2 to 22 seconds. The system is
covered by 547 automated backend tests and 85 frontend tests.
\end{abstract}

\tableofcontents
\newpage

% ==========================================================================
\section{Introduction}

\subsection{Background}

A digital twin is a software copy of a physical device, such as a car, a
machine or a sensor. Eclipse Ditto stores these twins, accepts updates from
devices and lets applications read and search them. A platform like this sits
in the critical path between a fleet of real devices and the applications that
depend on them.

Chaos engineering is the practice of deliberately breaking part of a system in
a controlled way, to learn how it fails and whether it recovers before a real
outage does the same thing unplanned.

\subsection{Problem statement}

A system that works on a healthy day is not necessarily resilient. Teams
usually discover how their platform fails during a real incident, when the
cost is highest. Existing chaos tools can inject faults, but injecting a fault
safely needs more: a way to stop automatically, proof that the fault was
removed, and a clear record of what was learned.

\subsection{Objectives}

\begin{enumerate}
  \item Inject controlled, reversible faults into a running Ditto deployment.
  \item Monitor system health continuously during an experiment.
  \item Stop and roll back automatically when a safety limit is crossed.
  \item Verify every rollback and measure recovery time.
  \item Turn measurements into findings and recommendations.
  \item Keep every safety decision in deterministic code, with AI as an
        optional advisory layer.
  \item Present everything in a live dashboard.
\end{enumerate}

\subsection{Scope}

The tool targets a local, disposable Ditto sandbox running under Docker
Compose. It is not intended for production systems. All device telemetry is
synthetic and labelled as such.

% ==========================================================================
\section{System Overview}

\subsection{System under test}

Eclipse Ditto runs as nine Docker containers: \texttt{gateway},
\texttt{things}, \texttt{things-search}, \texttt{policies},
\texttt{connectivity}, \texttt{mongodb}, \texttt{nginx}, and two user
interfaces. MongoDB is the shared dependency of the core services, which is
why the tool's dependency map rates it as the highest-risk service.

\subsection{Synthetic fleet}

A simulator creates digital twins and sends telemetry so that there is real
traffic to measure. Six product profiles are available.

\begin{table}[h]
\centering
\begin{tabularx}{\textwidth}{@{}lL@{}}
\toprule
\textbf{Profile} & \textbf{Example features} \\
\midrule
Electric vehicle & battery level, location, speed, temperature \\
Industrial machine & motor temperature, rpm, vibration, fault codes \\
Smart home & climate, thermostat, lighting, door contact \\
Energy system & solar generation, grid power, storage \\
Environmental sensor & temperature, humidity, air quality \\
Connected appliance & cycle status, power, error code \\
\bottomrule
\end{tabularx}
\caption{Product profiles in the simulated fleet.}
\end{table}

Every update carries a sequence number and a send time, so staleness and lost
updates are measured the same way for every product.

% ==========================================================================
\section{Architecture}

The system is organised in layers. Each layer calls only the one below it, and
exactly one module is allowed to start a real fault.

\begin{table}[h]
\centering
\begin{tabularx}{\textwidth}{@{}lLl@{}}
\toprule
\textbf{Layer} & \textbf{Responsibility} & \textbf{Main modules} \\
\midrule
Dashboard & Eight pages showing live state and controls & \texttt{frontend/} \\
API & REST endpoints and a live event stream & \texttt{api.py}, \texttt{events.py} \\
Analysis & Agents, planning, learning, recommendations & \texttt{agents.py}, \texttt{learning.py} \\
Control & The single approval gate; operating modes & \texttt{orchestrator.py}, \texttt{controller.py} \\
Fault engine & Fault types, scenarios, kill switch, rollback & \texttt{faults.py}, \texttt{scenario.py}, \texttt{guard.py} \\
Measurement & Health probe and baseline & \texttt{probe.py}, \texttt{monitor.py} \\
Traffic & Synthetic fleet & \texttt{simulator.py}, \texttt{profiles.py} \\
Storage & Knowledge store, records, recovery journal & \texttt{store.py} \\
Target & Eclipse Ditto in Docker & external \\
\bottomrule
\end{tabularx}
\caption{Architecture layers, top to bottom.}
\end{table}

\subsection{Design principles}

\begin{itemize}
  \item \textbf{One approval gate.} The command line and the dashboard both go
        through \texttt{orchestrator.py}, so there is one allowlist and one
        authorization check regardless of entry point.
  \item \textbf{Fail closed.} A missing, unhealthy or untrusted measurement is
        treated as a reason to stop, never as ``healthy''.
  \item \textbf{Verify, do not assume.} A rollback counts only after the clean
        state has been observed.
  \item \textbf{AI is advisory.} AI output is schema-validated before use and
        has a deterministic fallback.
\end{itemize}

% ==========================================================================
\section{Fault Injection}

Nine fault types are supported, in three groups.

\begin{table}[h]
\centering
\begin{tabularx}{\textwidth}{@{}llL@{}}
\toprule
\textbf{Group} & \textbf{Fault} & \textbf{Effect} \\
\midrule
Process & \texttt{kill} & Force-kills the container, then starts it again \\
        & \texttt{stop} & Graceful shutdown, then start \\
        & \texttt{restart} & Restarts the container \\
        & \texttt{pause} & Freezes the process; connections stay open \\
Network & \texttt{delay} & Fixed delay on outgoing packets \\
        & \texttt{jitter} & Variable delay per packet \\
        & \texttt{loss} & Drops a percentage of packets \\
        & \texttt{bandwidth} & Caps the outgoing rate \\
Resource & \texttt{cpu\_throttle} & Caps the container's CPU share \\
\bottomrule
\end{tabularx}
\caption{Supported fault types.}
\end{table}

Three further types (\texttt{cpu\_stress}, \texttt{memory\_pressure},
\texttt{request\_failure}) are deliberately not supported, because they cannot
be bounded or reversed safely in this deployment. The tool lists them together
with the reason.

Several faults can be combined in one scenario. A scenario is limited to three
concurrent faults across two services, and only combinations that have been
tested are allowed.

% ==========================================================================
\section{Safety Model}

\subsection{Kill switch}

A health probe runs every two seconds during an experiment. The experiment is
stopped and reverted when any of the following holds.

\begin{table}[h]
\centering
\begin{tabularx}{\textwidth}{@{}lL@{}}
\toprule
\textbf{Signal} & \textbf{Stop condition} \\
\midrule
Service outage & Direct reads or searches fail on three consecutive checks \\
Error rate & Above 5\% over the recent window \\
Latency & 95th percentile above 3$\times$ the healthy baseline (minimum 200 ms) \\
Staleness & Twin data more than 15 seconds old \\
Lost update & Any acknowledged update is missing; stops immediately \\
\bottomrule
\end{tabularx}
\caption{Kill-switch conditions (default thresholds).}
\end{table}

\subsection{Guarantees}

\begin{itemize}
  \item A real fault needs explicit operator authorization. An AI proposal can
        never run on its own.
  \item Every applied fault is rolled back independently and the result is
        verified.
  \item A revert is attempted on a kill-switch trip, at the time limit, on a
        manual stop and after a crash. A recovery journal on disk lets the
        tool finish a rollback after a restart.
  \item If a rollback cannot be verified, the system enters a
        \emph{recovery required} state and refuses all fault injection until
        an operator resolves it.
\end{itemize}

\subsection{Dry run}

A dry run walks the complete lifecycle with real health checks and lists every
command it would execute, but executes none. It needs no authorization and
never counts as evidence.

\subsection{Operating modes}

\begin{table}[h]
\centering
\begin{tabularx}{\textwidth}{@{}lL@{}}
\toprule
\textbf{Mode} & \textbf{What is allowed} \\
\midrule
\texttt{observe} & Read-only. Nothing can be injected. \\
\texttt{recommend} & Default. An operator may authorize a scenario the planner proposed. \\
\texttt{manual} & An operator may write and authorize any scenario that passes policy. \\
\texttt{autonomous\_sandbox} & A bounded loop runs experiments itself. Must be enabled in the configuration file; cannot be enabled from the dashboard. \\
\bottomrule
\end{tabularx}
\caption{Operating modes.}
\end{table}

% ==========================================================================
\section{Experiment Workflow}

\begin{enumerate}
  \item \textbf{Policy check.} Targets are allowlisted, parameters are within
        bounds and scenario limits are respected.
  \item \textbf{Preflight.} Docker and every target container are inspected,
        read-only.
  \item \textbf{Baseline.} Health is measured. If the saved baseline is older
        than 60 minutes, a fresh one is measured over 30 seconds. Nothing is
        injected unless Ditto is healthy.
  \item \textbf{Inject.} The fault is applied.
  \item \textbf{Monitor.} The probe runs every two seconds and the kill switch
        is evaluated on each result.
  \item \textbf{Revert and verify.} The fault is removed and the clean state
        is confirmed.
  \item \textbf{Recovery.} The time until three consecutive healthy checks is
        recorded.
  \item \textbf{Analysis.} Findings, a learning record and recommendations are
        produced. No recommendation is applied automatically.
\end{enumerate}

Each step writes to an activity log that records what the controller did and
why, including the measured numbers behind every continue-or-stop decision.

% ==========================================================================
\section{Analysis and Learning}

Seven agents process each experiment, coordinated by a single coordinator.

\begin{table}[h]
\centering
\begin{tabularx}{\textwidth}{@{}lL@{}}
\toprule
\textbf{Agent} & \textbf{Role} \\
\midrule
Risk & Rates each service from the dependency map and past runs \\
Fault injection & Proposes the next scenario from validated candidates \\
Recovery & Reports rollback and recovery per fault \\
Validation & Checks whether the run is usable as evidence \\
Insights & Derives findings, separating measured facts from hypotheses \\
Learning & Compares the outcome with the prediction and updates rankings \\
Remediation & Produces recommendations and tracks their status \\
\bottomrule
\end{tabularx}
\caption{Agents.}
\end{table}

Results are kept in a SQLite knowledge store. A deterministic ranking uses the
dependency map, service risk, experiment history and coverage gaps to suggest
which experiment is most worth running next. There is no trained model, by
design.

After an operator applies a fix, the original experiment can be re-run. The
tool then shows each metric before and after, and labels the result as an
observed correlation, not proof of causation.

\subsection{Role of AI}

When enabled, a language model (Google Gemini) may explain results, suggest
root causes and choose among candidates the policy engine has already
validated. It may not decide when to stop, authorize a run, issue a command or
apply a fix. With AI disabled, a rule-based analysis of the same measurements
is used and labelled as such.

% ==========================================================================
\section{Dashboard}

The dashboard is a single-page web application with eight pages.

\begin{table}[h]
\centering
\begin{tabularx}{\textwidth}{@{}lL@{}}
\toprule
\textbf{Page} & \textbf{Content} \\
\midrule
Overview & Target health against baseline, container states, alerts \\
Live monitoring & Latency, failure, staleness and consistency charts with fault intervals shaded and stop limits drawn \\
Products and twins & Fleet configuration and per-twin status \\
Experiment control & Operating mode, scenario builder, preview, dry run, authorization, emergency stop \\
Agents & State and last result of each agent \\
History & Past experiments with measurements, findings and timeline \\
Insights and learning & Evidence counts, rankings, recommendation tracking \\
Dependencies & Call edges, blast radius and risk per service \\
\bottomrule
\end{tabularx}
\caption{Dashboard pages.}
\end{table}

The interface uses a light theme by default with a dark option. Status is
never shown by colour alone: every state has a text label. All text colour
pairs meet a contrast ratio of at least 4.5:1.

% ==========================================================================
\section{Technology Stack}

\begin{table}[h]
\centering
\begin{tabularx}{\textwidth}{@{}lL@{}}
\toprule
\textbf{Area} & \textbf{Technology} \\
\midrule
Backend & Python 3.11, FastAPI, Uvicorn, Pydantic, httpx \\
Storage & SQLite, JSON experiment records \\
Frontend & React 18, TypeScript, Vite, Tailwind CSS, Recharts, React Router \\
Infrastructure & Docker Compose; Linux \texttt{tc netem} for network faults \\
Testing & pytest, Vitest, Playwright \\
AI (optional) & Google Gemini, with a deterministic fallback \\
\bottomrule
\end{tabularx}
\caption{Technology stack.}
\end{table}

% ==========================================================================
\section{Testing}

\begin{table}[h]
\centering
\begin{tabularx}{\textwidth}{@{}lL@{}}
\toprule
\textbf{Suite} & \textbf{Result} \\
\midrule
Backend unit and contract tests & 547 passed, 11 skipped (the skipped tests are opt-in suites that need a real Docker daemon or a live Ditto) \\
Frontend component tests & 85 passed \\
Type check and production build & Clean \\
Browser smoke test & All pages opened; plan preview, dry run, ranked candidates, experiment detail and re-run panel exercised; no console errors or failed requests \\
\bottomrule
\end{tabularx}
\caption{Test results.}
\end{table}

Backend tests mock Docker, Ditto and the AI provider, so they run offline.
They cover the fault registry, the scenario engine (partial failure, emergency
stop, lockout, journal replay), the agents, the learning logic and the HTTP
contract.

% ==========================================================================
\section{Results}

Six faults were injected one after another against the local sandbox, each
with a 15-second limit. All six were rolled back and verified, and Ditto
returned to a healthy state each time.

\begin{table}[h]
\centering
\begin{tabularx}{\textwidth}{@{}llLr@{}}
\toprule
\textbf{Fault} & \textbf{Target} & \textbf{What ended it} & \textbf{Recovery} \\
\midrule
\texttt{kill} & things-search & Kill switch: search failed 3 checks in a row & 22.1 s \\
\texttt{stop} & things-search & Kill switch: search failed 3 checks in a row & 12.1 s \\
\texttt{restart} & things-search & Ran its full 15 seconds & 2.0 s \\
\texttt{pause} & things-search & Kill switch: search failed 3 checks in a row & 4.0 s \\
\texttt{delay} (300 ms) & mongodb & Kill switch: latency above 3$\times$ baseline & 6.1 s \\
\texttt{jitter} (200$\pm$100 ms) & things & Kill switch: direct reads failed 3 checks in a row & 8.1 s \\
\bottomrule
\end{tabularx}
\caption{Live fault-injection run. Each experiment took 26 to 48 seconds end to end.}
\end{table}

\subsection{Observations}

\begin{itemize}
  \item The kill switch stopped five of the six experiments before their time
        limit, in each case for the signal the fault was expected to affect.
  \item A 300 ms delay on MongoDB was enough to push read latency past three
        times the baseline, which confirms MongoDB as the most sensitive
        dependency.
  \item Recovery after a forced kill took about ten times longer than after a
        restart (22.1 s against 2.0 s).
\end{itemize}

The remaining three fault types (\texttt{loss}, \texttt{bandwidth},
\texttt{cpu\_throttle}) were not part of this run.

% ==========================================================================
\section{Limitations}

\begin{itemize}
  \item The tool is built and tested for a local sandbox, not production.
  \item The automatic removal of a network fault survives a crash of the tool
        itself, but not a crash of the Docker daemon or the host.
  \item Staleness assumes one shared clock; it has not been validated across
        machines.
  \item The autonomous mode has been tested only against a simulated Docker.
  \item The largest fleet validated under load is 120 twins.
  \item AI analysis can be wrong about the mechanism even when the
        measurements are correct, so findings are a lead to verify, not a
        verdict.
\end{itemize}

% ==========================================================================
\section{Future Work}

\begin{itemize}
  \item Run the three remaining fault types and the autonomous loop against
        the live sandbox.
  \item Support deployments on Kubernetes in addition to Docker Compose.
  \item Add a fault-injecting proxy so that request-level failures can be
        tested safely.
  \item Validate larger fleets and multi-machine setups.
\end{itemize}

% ==========================================================================
\section{Conclusion}

ditto-chaos shows that chaos experiments on a digital-twin platform can be run
with strong safety guarantees: an automatic kill switch, verified rollback, a
lockout when recovery cannot be confirmed, and explicit operator
authorization. Keeping every safety decision in deterministic code, with AI
limited to explanation and ranking, made the system both testable and
trustworthy. In a live run the tool injected six different faults, stopped
five of them automatically, and verified recovery in every case.

% ==========================================================================
\section*{References}
\addcontentsline{toc}{section}{References}

\begin{enumerate}[label={[\arabic*]}]
  \item Eclipse Ditto. \url{https://github.com/eclipse-ditto/ditto}
  \item ditto-chaos source code. \url{https://github.com/NJ800/ditto-chaos}
  \item FastAPI. \url{https://fastapi.tiangolo.com}
  \item React. \url{https://react.dev}
  \item Docker Compose. \url{https://docs.docker.com/compose/}
\end{enumerate}

\end{document}
